% Part: incompleteness
% Chapter: representability-in-q
% Section: comp-representable

\documentclass[../../../include/open-logic-section]{subfiles}

\begin{document}

\olfileid{inc}{req}{crq}
\olsection{संगणनक्षम फलने~$\Th{Q}$ मध्ये निरूपणीय असतात}

\begin{thm}
प्रत्येक संगणनक्षम फलन~$\Th{Q}$ मध्ये निरूपणीय असते.
\end{thm}

\begin{proof}
संदिग्धता टाळण्यासाठी, चर्च--ट्यूरिंग प्रबंध वापरून फलन
संगणनक्षम असते तेव्हाच आणि तरच ते सामान्य पुनरावर्ती असते,
असे म्हणू. सामान्य पुनरावर्ती फलने म्हणजे शून्य फलन~$\Zero$,
उत्तरवर्ती फलन~$\Succ$ आणि प्रक्षेपण फलन~$\Proj{n}{i}$
यांपासून संयोजन, आदिम पुनरावर्तन व नियमित लघुतमीकरण वापरून
परिभाषित करता येणारी फलने. \olref[pri]{lem:prim-rec} नुसार
$f$ आणि~$g$ यांपासून आदिम पुनरावर्तनाने परिभाषित होणारे
कोणतेही फलन~$h$ हे $f$, $g$, $\Zero$,
$\Succ$, $\Proj{n}{i}$, $\Add$, $\Mult$, $\Char{=}$
यांपासून संयोजन व नियमित लघुतमीकरण वापरूनही परिभाषित
करता येते. परिणामी, फलन सामान्य पुनरावर्ती असते तेव्हाच
आणि तरच ते $\Zero$, $\Succ$, $\Proj{n}{i}$, $\Add$,
$\Mult$, $\Char{=}$ यांपासून संयोजन व नियमित लघुतमीकरण
वापरून परिभाषित करता येते.

शिवाय, ही सर्व मूलभूत फलने~$\Th{Q}$ मध्ये निरूपणीय आहेत,
हे आपण सिद्ध केले आहे
(\cref{inc:req:bre:prop:rep-zero,inc:req:bre:prop:rep-succ,inc:req:bre:prop:rep-proj,inc:req:bre:prop:rep-id,inc:req:bre:prop:rep-add,inc:req:bre:prop:rep-mult});
आणि निरूपणीय फलनांपासून संयोजन किंवा नियमित लघुतमीकरणाने
परिभाषित होणारे कोणतेही फलन
(\olref[cmp]{prop:rep-composition},
\olref[min]{prop:rep-minimization}) हेही निरूपणीय असते.
म्हणून प्रत्येक सामान्य पुनरावर्ती फलन~$\Th{Q}$ मध्ये
निरूपणीय असते.
\end{proof}

\begin{explain}
संगणनक्षम फलनांचा वर्ग म्हणजेच $\Th{Q}$ मध्ये निरूपणीय
फलनांचा वर्ग, असे लक्षणन आपल्याला मिळाले आहे. प्रत्यक्षात
हा पुरावा अधिक व्यापक आहे. निरूपणीयतेच्या व्याख्येवरून
$\Th{Q}$ चा विस्तार असलेल्या कोणत्याही उपपत्तीत (किंवा
ज्यात $\Th{Q}$ चे औपचारिक अर्थनिर्धारण करता येते अशा उपपत्तीत)
संगणनक्षम फलनांचे निरूपण करता येते, हे सहज दिसते. उलट,
ज्या सुसंगत !!{derivation} व्यवस्थेत भिन्न संख्यांकांची असमता
सिद्ध होते आणि !!{derivation} ओळखण्याची प्रक्रिया संगणनक्षम
असते, त्या व्यवस्थेत निरूपणीय असलेले प्रत्येक फलन संगणनक्षम
असते. विसंगत व्यवस्थेत प्रत्येक सूत्र सिद्ध होत असल्याने
ही उलटी दिशा तेथे लागू पडत नाही. म्हणून, उदाहरणार्थ, संगणनक्षम
फलनांचा वर्ग म्हणजे पेआनो अंकगणितात किंवा अगदी
झर्मेलो--फ्रेंकेल संच उपपत्तीत निरूपणीय फलनांचा वर्ग,
असेही लक्षणन करता येते. गोडेल यांनी नोंदवल्याप्रमाणे,
हे काहीसे आश्चर्यकारक आहे. सिद्धतायोग्यतेबाबत प्रश्नांचे
उत्तर कोणती उपपत्ती निवडली आहे यावर फार अवलंबून असते,
हे आपण पुढे पाहू; साधारणपणे स्वयंसिद्धके जितकी बलवान,
तितके अधिक सिद्ध करता येते. पण स्वयंसिद्धकीय उपपत्तींच्या
मोठ्या वर्गात निरूपणीय फलने नेमकी संगणनक्षम फलनेच असतात;
त्या उपपत्ती सुसंगत आणि निर्णेय स्वयंसिद्धकसंचाने मांडता येत असतील,
तर अधिक बलवान उपपत्ती अधिक फलनांचे निरूपण करत नाहीत.
\end{explain}

\end{document}
